
\pagebreak

\chapter{React basics.}


\medskip

React.js is a powerful JavaScript library for building user interfaces. It allows developers to create dynamic and responsive web applications efficiently. One of its key features is its declarative and component-based architecture, which makes UI development more structured and manageable.

React was developed by Jordan Walke, a software engineer at Facebook, and was first released in 2013. It was created to improve Facebook's complex UI development by introducing a more efficient way to manage updates and rendering. Over time, React gained widespread adoption due to its performance benefits and reusability of components.

React development is streamlined with automatic reload and hot reloading features. When changes are made to the code, they are instantly reflected in the browser without needing a full page refresh. This significantly speeds up the development process and improves the developer experience.

React facilitates the development of Single Page Applications (SPA). Unlike traditional multi-page applications, an SPA loads a single HTML file and dynamically updates the content as needed without requiring a page reload. This results in faster navigation and a more seamless user experience.

It is important to understand difference between React and React Native. Whereas React is used for building web applications, React Native is a framework used primarily for mobile apps. React runs in browsers, React Native on Native UI components on iOS and Android.

There were several major versions of React:

\begin{itemize}
\item React 16 (2017): Introduced the Fiber reconciliation engine for better performance.

\item React 17 (2020): Improved gradual upgrades without breaking changes.

\item React 18 (2022): Added features like Concurrent Rendering and Suspense for better performance.

\item Latest Version: React continues to evolve with new features and optimizations.
\end{itemize}

\medskip

\textbf{Example simple React application}

\medskip

Let's use create-react-app tool to create a new react app:

\begin{lstlisting}
npx create-react-app new-app
\end{lstlisting}

After that:

\begin{lstlisting}
cd new-app
\end{lstlisting}

And:

\begin{lstlisting}
npm start
\end{lstlisting}

React development is streamlined with automatic reload and hot reloading features. When changes are made to the code, they are instantly reflected in the browser without needing a full page refresh. This significantly speeds up the development process and improves the developer experience.

Lets build single page app, that returns single html file. Take a look at a generated repository.
The index.js file serves as the entry point of a React application. It is responsible for importing necessary modules and rendering the root component into the DOM.
The ReactDOM.render() function (or its modern equivalent, createRoot) is executed to render the main App component inside a specific DOM element, usually an element with the id root in index.html.

The App.js file contains the main application component. In React, components are written as JavaScript functions that return JSX. JSX (JavaScript XML) allows developers to write HTML-like code within JavaScript, making UI development more intuitive.


Let's rewrite the App.js component:

\begin{lstlisting}
import './App.css';
import React from 'react';

function App() {
  return (
    <div className='container'>
      <h2>Student list </h2>
      <ul className='stud-list'>
        <li>Ivanov</li>
        <li>Petrov</li>
        <li>Frolov</li>
      </ul>
    </div>
  );
}

export default App;
\end{lstlisting}

Note: we can set styles in App.css, for example, class container.
As the word class is reserved in JS, instead of class of the element we write className.

Add the following in App.css file:

\begin{lstlisting}
.container {
    max-width: 400px;
    margin: 20px auto;
    padding: 20px;
    background: #f8f9fa;
    border-radius: 10px;
    box-shadow: 0 4px 6px rgba(0, 0, 0, 0.1);
    text-align: center;
}

h2 {
    font-size: 1.5rem;
    color: #333;
    margin-bottom: 15px;
}

.stud-list {
    list-style: none;
    padding: 0;
}

    .stud-list li {
        background: #ffffff;
        padding: 10px;
        margin: 5px 0;
        border-radius: 5px;
        box-shadow: 0 2px 4px rgba(0, 0, 0, 0.1);
        transition: background 0.3s ease-in-out;
    }

        .stud-list li:hover {
            background: #e3e6ea;
            cursor: pointer;
        }
\end{lstlisting}

List element should be outsourced to the components.
Add the components folder in \verb"src", there add StudList.js file.

There we add the following:

\begin{lstlisting}
import React from 'react';

const StudList = () => {
    return (
        <ul className='stud-list'>
        <li>Ivanov</li>
        <li>Petrov</li>
        <li>Frolov</li>
      </ul>
    )
}

export default StudList;
\end{lstlisting}

And change the App.js to the following:

\begin{lstlisting}
import './App.css';
import React from 'react';
import StudList from './components/StudList';

function App() {
  return (
    <div className='container'>
      <h2>Student list </h2>
      <StudList />
    </div>
  );
}

export default App;
\end{lstlisting}

If we save now, we see the same output and the same styling as before.
But it is still a good practice to add a separate css file for a component.
StudList.css in components folder:

\begin{lstlisting}
.stud-list {
    list-style: none;
    padding: 0;
}

    .stud-list li {
        background: #ffffff;
        padding: 10px;
        margin: 5px 0;
        border-radius: 5px;
        box-shadow: 0 2px 4px rgba(0, 0, 0, 0.1);
        transition: background 0.3s ease-in-out;
    }

        .stud-list li:hover {
            background: #e3e6ea;
            cursor: pointer;
        }
\end{lstlisting}

\medskip
\textbf{React Props.}

With splitting our code into components, we have now achieved separation of concerns. But how do we pass data between the components?
Here useful concept to introduce is the concept of props. Props are arguments that are passed automatically by react to the component.

Let's rewrite \verb"App.js":

\begin{lstlisting}
import './App.css';
import React from 'react';
import StudList from './components/StudList';

function App() {
  const sutdList = [
    {id:1, text: "Ivanov"},
    {id:2, text: "Petrov"},
    {id:3, text: "Frolov"}
  ];

  return (
    <div className='container'>
      <h2>Student list </h2>
      <StudList studs={sutdList}/>
    </div>
  );
}

export default App;
\end{lstlisting}

\verb"StudList.js": let's map array of elements to an array of jsx objects. Also added \verb"console.log" to see, what comes in props.

\begin{lstlisting}
import React from 'react';
import './StudList.css'

const StudList = props => {
    console.log(props.studs);

    return (
        <ul className='stud-list'>
            {
                props.studs.map((stud) => {
                    return <li>{stud.text}</li>
                })
            }
        </ul>
    )
}

export default StudList;
\end{lstlisting}

\medskip
\textbf{React Events}

Add component NewStudent.js

\begin{lstlisting}
import React from 'react';
import './NewStudent.css'

const NewStudent = props => {
    const addStudentHandler = event => {
        event.preventDefault();

        const newStud = {
            id: 1,
            text: 'text'
        }

        props.onAddStud(newStud);
    }

    return (
    <form className='new-student' onSubmit={addStudentHandler}>
        <input type='text'/>
        <button type='submit'>Add Student</button>
    </form>
    )
}

export default NewStudent;
\end{lstlisting}

Add styles for it:

\begin{lstlisting}
.new-student {
  display: flex;
  gap: 10px;
  max-width: 400px;
  margin: 20px auto;
  padding: 15px;
  background: #f8f9fa;
  border-radius: 10px;
  box-shadow: 0 4px 6px rgba(0, 0, 0, 0.1);
}

.new-student input {
  flex: 1;
  padding: 10px;
  border: 2px solid #ccc;
  border-radius: 5px;
  font-size: 1rem;
  outline: none;
  transition: border-color 0.3s;
}

.new-student input:focus {
  border-color: #007bff;
}

.new-student button {
  padding: 10px 15px;
  background: #007bff;
  color: #fff;
  border: none;
  border-radius: 5px;
  cursor: pointer;
  font-size: 1rem;
  transition: background 0.3s ease-in-out;
}

.new-student button:hover {
  background: #0056b3;
}

\end{lstlisting}

And in \verb"App.js":

\begin{lstlisting}
import './App.css';
import React from 'react';
import StudList from './components/StudList';
import NewStudent from './components/NewStudent/NewStudent'

function App() {
  const studList = [
    {id:1, text: "Ivanov"},
    {id:2, text: "Petrov"},
    {id:3, text: "Frolov"}
  ];

  const addNewStudHandler = (newStud) => {
    studList.push(newStud);
  }

  return (
    <div className='container'>
      <h2>Student list </h2>
      <NewStudent onAddStud={addNewStudHandler} />
      <StudList studs={studList}/>
    </div>
  );
}

export default App;
\end{lstlisting}

Now every time \verb"addNewStudHandler" method is triggered in the child component via submittion of the form, the parent event is notified and has a new entry in the array.


\medskip
\textbf{React State}

State in React is a way to manage dynamic data within a component. It allows components to store, update, and react to changes in data over time. State is typically declared using the useState hook in functional components.

In App.js:

\begin{lstlisting}
import './App.css';
import React, {useState} from 'react';
import StudList from './components/StudList';
import NewStudent from './components/NewStudent/NewStudent';


function App() {
  const [studList, setStud ] = useState(
    [
      {id:1, text: "Ivanov"},
      {id:2, text: "Petrov"},
      {id:3, text: "Frolov"}
    ]
  );

  // const studList = [
  //   {id:1, text: "Ivanov"},
  //   {id:2, text: "Petrov"},
  //   {id:3, text: "Frolov"}
  // ];

  const addNewStudHandler = (newStud) => {
    // setStud(studList.concat(newStud));
    setStud( (oldStuList) => {return oldStuList.concat(newStud)} );
  }

  return (
    <div className='container'>
      <h2>Student list </h2>
      <NewStudent onAddStud={addNewStudHandler} />
      <StudList studs={studList}/>
    </div>
  );
}

export default App;
\end{lstlisting}

Here setStud - method that allows to handle state updates. State can be anything. UseState always returns array with two elements - state and method to update it.



\medskip
\textbf{React Data binding}

Data binding in React is the process of linking data between the component's state or props and the UI. React primarily uses one-way data binding, meaning that data flows in a single direction�from the component's state or props to the UI. This makes it predictable and easier to debug.

In NewStud.js:

\begin{lstlisting}
import React, {useState} from 'react';
import './NewStudent.css'

const NewStudent = props => {
    const [enteredText, setText] = useState('');

    const addStudentHandler = event => {
        event.preventDefault();

        const newStud = {
            id: 1,
            text: enteredText
        }

        setText('');

        props.onAddStud(newStud);
    }

    const textChangeHandler = event => {
        setText(event.target.value);
    }

    return (
    <form className='new-student' onSubmit={addStudentHandler}>
        <input type='text' value={enteredText} onChange={textChangeHandler}/>
        <button type='submit'>Add Student</button>
    </form>
    )
}

export default NewStudent;
\end{lstlisting}


